\documentclass[11pt,twocolumn]{article}
\usepackage[a4paper,left=17mm,right=17mm,top=17mm,bottom=18mm,columnsep=7mm]{geometry}
\usepackage{amsmath,amssymb,mathtools,microtype}
\usepackage[hidelinks]{hyperref}
\usepackage{flushend}
\newcommand{\Q}{\mathcal Q}\newcommand{\W}{\mathcal W}
\newcommand{\sxdeg}{\operatorname{sxdeg}}
\setlength{\parindent}{1em}
\setlength{\abovedisplayskip}{6pt}\setlength{\belowdisplayskip}{6pt}
\allowdisplaybreaks
\begin{document}
\twocolumn[{\centering
{\Large\bfseries Infinite NPA depth, finite semidefinite lift\par}
\vspace{5pt}
{\normalsize An exact quantum region in the simplest Bell experiment\par}
\vspace{13pt}}]

The first NPA level recovers the quantum CHSH bound. Once marginal averages
are retained, every fixed level misses some quantum bounds. Yet the whole
two-dimensional region tested by symmetric doubly tilted CHSH has a fixed,
exact semidefinite representation. The distinction can be seen directly in
the proofs: equal measurement overlaps determine all its supporting lines,
while finite NPA certificates must hold for two independently varying angles.

Alice and Bob choose between two measurements with outcomes $\pm1$.
Write $a_x=\langle A_x\rangle$, $b_y=\langle B_y\rangle$ and
$E_{xy}=\langle A_xB_y\rangle$. Retain
\begin{equation}
 M=a_0+b_0,\qquad S=E_{00}+E_{01}+E_{10}-E_{11}.
 \label{eq:coordinates}
\end{equation}
Their quantum region $\W=\{(M,S):p\in\Q_{222}\}$ includes all quantum
behaviors, with no equality imposed on the marginal averages.
If $\W_L$ denotes its level-$L$ NPA outer approximation, then
\begin{equation}
 \boxed{\begin{gathered}
 \W\subsetneq\W_L\quad\text{for every finite }L,\\
 \sxdeg(\W)=2.
 \end{gathered}}
 \label{eq:result}
\end{equation}
We construct one lift with positive semidefinite (PSD) blocks of orders
five and four. Extension degree is the smallest possible largest block
order among all finite lifts; its value two follows separately from a
planar representation theorem.

\emph{Why one angle suffices.}
To determine $\W$, consider an arbitrary supporting inequality
$hS+\alpha M\leq\lambda$, with real $h,\alpha$.
Jordan's lemma reduces its validity to positivity of two-qubit slacks
$\lambda I-H$, where
\[
\begin{gathered}
 H=h\,\mathrm{CHSH}+\alpha(ZI+IZ),\\
 A_1=cZ+\sqrt{1-c^2}X,\quad B_1=dZ+\sqrt{1-d^2}X.
\end{gathered}
\]
Here $A_0=B_0=Z$ and $c,d\in[-1,1]$. Deterministic strategies attain
$\ell=2|h|+2|\alpha|$, so $\lambda\geq\ell$ is necessary.
Above this bound, the first three elementary symmetric functions of
the eigenvalues of $\lambda I-H$ are nonnegative. Its determinant
therefore decides positivity.

Put $x=(1-c)/2$, $y=(1-d)/2$, and let
$D(x,y)=\det(\lambda I-H)$. The decisive identity is
\begin{equation}
\begin{split}
 D(x,y)-D(\sqrt{xy},\sqrt{xy})
 ={}&64h^2(4h^2-\alpha^2)xy\\
 &{}\times(\sqrt x-\sqrt y)^2 .
\end{split}
\label{eq:key}
\end{equation}
In the nonclassical range $|\alpha|\leq|h|$, replacing the two overlaps
by equal ones preserves $xy$ and can only decrease the determinant.
Thus positivity at equal overlaps implies positivity at every pair.
For $|\alpha|\geq|h|$, the quantum support is the deterministic value
$\ell$: average the two signed one-sided tilted-CHSH bounds and add
the remaining marginal slack. Equal overlaps attain this value too.
Hence they suffice for \emph{every} support direction. Symmetry of the
functional alone would not give this conclusion.

At equal overlaps $c=d$, the singlet decouples. In the remaining
three-dimensional block, set $c=(1-t^2)/(1+t^2)$ and multiply the slack
by $(1+t^2)^2$. With
\[
 p_0=(1+t^2)^2,\qquad p_1=1-t^4,\qquad p_2=1+4t^2-t^4,
\]
the result is the quartic matrix
\begin{equation}
 K(t)=
 \begin{pmatrix}
 \lambda p_0-2hp_1&-2\alpha p_0&0\\
 -2\alpha p_0&\lambda p_0-2hp_2&-8ht^3\\
 0&-8ht^3&\lambda p_0+2hp_2
 \end{pmatrix}.
 \label{eq:K}
\end{equation}
Its value at zero and leading coefficient force $\lambda\geq\ell$,
which also controls the singlet. The missing endpoint follows by
continuity. We have therefore proved
\begin{equation}
 \max_{\W}(hS+\alpha M)\leq\lambda
 \quad\Longleftrightarrow\quad K(t)\succeq0\quad\forall t\in\mathbb R .
 \label{eq:support}
\end{equation}
All three support coefficients still enter \emph{affinely}. This is
what will give one lift for the whole region.

\emph{From one variable to one finite lift.}
A PSD real univariate matrix polynomial is a polynomial matrix sum
of squares, including on its singular boundary~\cite{hs}.
For a quartic, the factors have degree at most two.
The parity of Eq.~\eqref{eq:K} splits the nine quadratic features into
\[
\begin{aligned}
 Z_+&=(\xi_1,\xi_2,t\xi_3,t^2\xi_1,t^2\xi_2)^T,\\
 Z_-&=(\xi_3,t\xi_1,t\xi_2,t^2\xi_3)^T .
\end{aligned}
\]
Consequently Eq.~\eqref{eq:support} is equivalent to
\begin{equation}
 \begin{gathered}
 \xi^TK(t)\xi=Z_+^TGZ_++Z_-^TJZ_-,\\
 G\succeq0,\quad J\succeq0,
 \end{gathered}
 \label{eq:gram}
\end{equation}
with $G$ of order five and $J$ of order four. Matching coefficients
gives sixteen affine equations. There is no remaining angle
optimization.

Dualizing these equations gives an exact membership test with the
same two block orders, sixteen auxiliary variables, and three affine
equations for normalization, $M$ and $S$. The explicit matrices are
given in the accompanying complete proofs.
\newpage
A positive definite Gram
pair at $(\lambda,\alpha,h)=(1,0,0)$ supplies strict feasibility;
its weighted-trace normalization bounds every dual variable. The
dual image is therefore closed and attains every boundary point.
Since its supporting half-spaces are exactly Eq.~\eqref{eq:support},
it equals $\W$.

\emph{Why finite NPA depth still fails.}
For the same region, let
\[
 \begin{gathered}
 \omega_{\rm Q}(\alpha)=\max_{\W}(S+\alpha M),\\
 \Delta_\alpha=\omega_{\rm Q}(\alpha)-(2+2\alpha).
 \end{gathered}
\]
The nontermination proof~\cite{prior} gives, with $T=1-\alpha$,
\begin{equation}
\begin{gathered}
 \Delta_{1-T}=\tfrac43T^3+O(T^4),\\
 \frac{\omega_{\rm NPA}^{(L)}(1-T)-\omega_{\rm Q}(1-T)}
 {\Delta_{1-T}}\longrightarrow+\infty
 \qquad(T\downarrow0)
\end{gathered}
\label{eq:npa}
\end{equation}
for each fixed finite $L$.

The obstruction survives in the two independent measurement angles.
Set $T=t^2$, use angles $tu,tv$, and the normalized version of
\[
 \xi_t=(1,t^3u^2v/12,t^3uv^2/12,-t^2uv/4)^T .
\]
The gap between the quantum bound and this strategy, divided by
$\Delta_{1-t^2}$, tends to $\mathcal M(u/2,v/2)$, where
\[
 {\cal M}(U,V)=1+U^4V^2+U^2V^4-3U^2V^2.
\]
A bounded ratio in Eq.~\eqref{eq:npa} would make the fixed-level
NPA square factors converge, after rescaling, to polynomial squares
representing $\mathcal M+c$ for some $c\geq0$.
The fixed word length and bounded Gram trace control their Taylor
remainders; bounded polynomial values then control their coefficients.
These compactness steps are proved in full in the companion.

The contradiction itself takes one line of coefficients.
Each square factor has degree at most three and is constant on both
axes, hence has the form $a+bUV+dU^2V+eUV^2$.
The coefficient of $U^2V^2$ in a sum of such squares is $\sum b^2$;
in $\mathcal M+c$ it is $-3$.

Equal overlaps thus determine the support values, while an NPA
certificate must represent the bound for independently chosen
measurement angles. The univariate squares in Eq.~\eqref{eq:gram}
are obtained after the support reduction. They need not be identities
in a fixed list of observable words.

\emph{Minimal block order and thirty-two exact projections.}
The set $\W$ is closed, convex and semialgebraic, so Scheiderer's
planar theorem gives $\sxdeg(\W)\leq2$~\cite{soc}.
The cubic term in Eq.~\eqref{eq:npa} excludes a polygon: a polygon's
support would be affine sufficiently close to $\alpha=1$.
Thus $\sxdeg(\W)=2$. The blocks $5,4$ give the explicit construction;
the planar theorem gives the optimal order.

Physical input and output relabellings give the same quantum region for
\[
 \left(a_x+\chi_{xy}b_y,\ \sum_{ij}\chi_{ij}E_{ij}\right),
 \qquad \prod_{ij}\chi_{ij}=-1.
\]
The eight CHSH sign patterns and four selected input pairs give
32 distinct planes. Intersecting their fixed lifts with any NPA
relaxation makes every support value on these planes exact, including
all the directions in Eq.~\eqref{eq:npa}.

\begin{thebibliography}{3}\setlength{\itemsep}{0pt}
\bibitem{prior} A. Chaturvedi, \emph{No Finite NPA Level Characterizes
the Complete Quantum Set in the Simplest Bell Scenario}, project
manuscript (September 2026). Its complete nontermination proof is
reproduced in the companion, Part II.
\bibitem{hs} C. Hanselka and R. Sinn, \emph{Math. Z.} \textbf{292},
83--101 (2019), Corollary 2.8;
\href{https://arxiv.org/abs/1707.08261}{arXiv:1707.08261}.
\bibitem{soc} C. Scheiderer, \emph{SIAM J. Appl. Algebra Geom.}
\textbf{5}, 114--139 (2021), Theorem 0.1;
\href{https://arxiv.org/abs/2004.04196}{arXiv:2004.04196}.
\end{thebibliography}
\end{document}
